\documentclass[11pt]{article}
\newcommand{\doctitle}{Annotated Bibliography: Monte Carlo
  Simulation Methods for N-of-1 Clinical Trials}
\newcommand{\shorttitle}{Annotated Bibliography}
\newcommand{\docdate}{2026-04-08}
% pmsimstats-ng standard document preamble
% Usage: % pmsimstats-ng standard document preamble
% Usage: % pmsimstats-ng standard document preamble
% Usage: \input{pmsimstats-preamble} after \documentclass[11pt]{article}
%
% Documents must define before \input:
%   \newcommand{\doctitle}{...}       Full document title
%   \newcommand{\shorttitle}{...}     Short title for header
%   \newcommand{\docdate}{YYYY-MM-DD} Document date

\usepackage[margin=1in]{geometry}
\usepackage{amsmath,amssymb}
\usepackage[T1]{fontenc}
\usepackage{lmodern}
\usepackage{microtype}
\usepackage{graphicx}
\graphicspath{{figures/}}
\usepackage{booktabs}
\usepackage{longtable}
\usepackage{array}
\usepackage{hyperref}
\usepackage{xcolor}
\usepackage{fancyhdr}
\usepackage{parskip}
\usepackage{caption}
\usepackage{enumitem}
\usepackage{verbatim}
\usepackage{listings}

\lstset{
  language=R,
  basicstyle=\small\ttfamily,
  keywordstyle=\color{blue!70!black},
  commentstyle=\color{green!50!black},
  stringstyle=\color{red!60!black},
  breaklines=true,
  frame=single,
  framerule=0.5pt,
  rulecolor=\color{gray!40},
  backgroundcolor=\color{gray!5},
  xleftmargin=1em,
  framexleftmargin=1em,
  columns=flexible
}

\hypersetup{
  colorlinks=true,
  linkcolor=blue!60!black,
  citecolor=green!50!black,
  urlcolor=blue!60!black
}

\pagestyle{fancy}
\fancyhf{}
\lhead{\footnotesize pmsimstats-ng Technical Report}
\rhead{\footnotesize \shorttitle}
\rfoot{\footnotesize \thepage}
\renewcommand{\headrulewidth}{0.4pt}

\title{\doctitle}
\author{pmsimstats team}
\date{\docdate}
 after \documentclass[11pt]{article}
%
% Documents must define before \input:
%   \newcommand{\doctitle}{...}       Full document title
%   \newcommand{\shorttitle}{...}     Short title for header
%   \newcommand{\docdate}{YYYY-MM-DD} Document date

\usepackage[margin=1in]{geometry}
\usepackage{amsmath,amssymb}
\usepackage[T1]{fontenc}
\usepackage{lmodern}
\usepackage{microtype}
\usepackage{graphicx}
\graphicspath{{figures/}}
\usepackage{booktabs}
\usepackage{longtable}
\usepackage{array}
\usepackage{hyperref}
\usepackage{xcolor}
\usepackage{fancyhdr}
\usepackage{parskip}
\usepackage{caption}
\usepackage{enumitem}
\usepackage{verbatim}
\usepackage{listings}

\lstset{
  language=R,
  basicstyle=\small\ttfamily,
  keywordstyle=\color{blue!70!black},
  commentstyle=\color{green!50!black},
  stringstyle=\color{red!60!black},
  breaklines=true,
  frame=single,
  framerule=0.5pt,
  rulecolor=\color{gray!40},
  backgroundcolor=\color{gray!5},
  xleftmargin=1em,
  framexleftmargin=1em,
  columns=flexible
}

\hypersetup{
  colorlinks=true,
  linkcolor=blue!60!black,
  citecolor=green!50!black,
  urlcolor=blue!60!black
}

\pagestyle{fancy}
\fancyhf{}
\lhead{\footnotesize pmsimstats-ng Technical Report}
\rhead{\footnotesize \shorttitle}
\rfoot{\footnotesize \thepage}
\renewcommand{\headrulewidth}{0.4pt}

\title{\doctitle}
\author{pmsimstats team}
\date{\docdate}
 after \documentclass[11pt]{article}
%
% Documents must define before \input:
%   \newcommand{\doctitle}{...}       Full document title
%   \newcommand{\shorttitle}{...}     Short title for header
%   \newcommand{\docdate}{YYYY-MM-DD} Document date

\usepackage[margin=1in]{geometry}
\usepackage{amsmath,amssymb}
\usepackage[T1]{fontenc}
\usepackage{lmodern}
\usepackage{microtype}
\usepackage{graphicx}
\graphicspath{{figures/}}
\usepackage{booktabs}
\usepackage{longtable}
\usepackage{array}
\usepackage{hyperref}
\usepackage{xcolor}
\usepackage{fancyhdr}
\usepackage{parskip}
\usepackage{caption}
\usepackage{enumitem}
\usepackage{verbatim}
\usepackage{listings}

\lstset{
  language=R,
  basicstyle=\small\ttfamily,
  keywordstyle=\color{blue!70!black},
  commentstyle=\color{green!50!black},
  stringstyle=\color{red!60!black},
  breaklines=true,
  frame=single,
  framerule=0.5pt,
  rulecolor=\color{gray!40},
  backgroundcolor=\color{gray!5},
  xleftmargin=1em,
  framexleftmargin=1em,
  columns=flexible
}

\hypersetup{
  colorlinks=true,
  linkcolor=blue!60!black,
  citecolor=green!50!black,
  urlcolor=blue!60!black
}

\pagestyle{fancy}
\fancyhf{}
\lhead{\footnotesize pmsimstats-ng Technical Report}
\rhead{\footnotesize \shorttitle}
\rfoot{\footnotesize \thepage}
\renewcommand{\headrulewidth}{0.4pt}

\title{\doctitle}
\author{pmsimstats team}
\date{\docdate}


\begin{document}
\maketitle
\thispagestyle{fancy}

This annotated bibliography surveys the simulation
methodology literature relevant to the design and analysis
of aggregated N-of-1 clinical trials. The emphasis falls on
Monte Carlo simulation approaches that inform power
analysis, carryover modeling, biomarker-treatment
interaction estimation, and Bayesian inference in crossover
designs. The motivating clinical application is the
prazosin-PTSD trial program.

% --------------------------------------------------------
\section{Core Simulation Frameworks}
% --------------------------------------------------------

\textbf{Thomas RG, Hendrickson RC, Schork NJ, Raskind MA.}
Optimizing aggregated N-of-1 trial designs for predictive
biomarker validation: statistical methods and theoretical
findings. \textit{Frontiers in Digital Health}.
2020;2:13. doi:10.3389/fdgth.2020.00013.

\smallskip\noindent
This paper develops the simulation framework that motivates
the present project. The authors derive analytic power
formulae for aggregated N-of-1 trials designed to detect
predictive biomarker effects and validate them through
Monte Carlo simulation. Key design factors include the
number of treatment cycles, the number of participants,
the magnitude of the biomarker-treatment interaction, and
within-subject serial correlation. The principal finding is
that aggregated N-of-1 designs can achieve adequate power
to detect moderate biomarker-treatment interactions with
substantially fewer participants than parallel-group
designs.

\medskip
\textbf{Blackston JW, Chapple AG, McGree JM, McDonald S,
Nikles J.} Comparison of aggregated N-of-1 trials with
parallel and crossover randomized controlled trials using
simulation studies. \textit{Healthcare}. 2019;7(4):137.
doi:10.3390/healthcare7040137.

\smallskip\noindent
Blackston et al.\ conduct 5000 simulated trials per
scenario across four variance structures to compare
aggregated N-of-1, parallel RCT, and crossover designs.
The simulation varies patient heterogeneity, model error,
and carryover strength in three-cycle designs. Aggregated
N-of-1 trials consistently outperform both traditional
designs in power and required sample size, though they
exhibit elevated type~I error under moderate-to-strong
unmodeled carryover. The study reinforces the need to
explicitly model carryover when planning N-of-1 trials.

\medskip
\textbf{Wang Y, Schork NJ.} Power and design issues in
crossover-based N-of-1 clinical trials with fixed data
collection periods. \textit{Healthcare}. 2019;7(3):84.
doi:10.3390/healthcare7030084.

\smallskip\noindent
Wang and Schork use analytical power formulae under an
AR(1) serial correlation model, validated by simulation, to
examine the effects of serial correlation,
heteroscedasticity, washout periods, and carryover on
N-of-1 trial power. The investigation systematically varies
the number of evaluation periods and the autocorrelation
coefficient. Results demonstrate that serial correlation
can substantially reduce power, that increasing the number
of evaluation periods partially compensates for this loss,
and that certain carryover effects are difficult to detect
without careful design choices.

% --------------------------------------------------------
\section{Power and Sample Size}
% --------------------------------------------------------

\textbf{Senn S.} Sample size considerations for N-of-1
trials. \textit{Statistical Methods in Medical Research}.
2019;28(2):372--383. doi:10.1177/0962280217726801.

\smallskip\noindent
Senn identifies five distinct criteria for planning sample
sizes in N-of-1 trials and provides formulae for each under
a normal-normal mixture model with variance components for
within-cycle variation and treatment-by-patient interaction.
The design addressed is randomization to treatment and
control within paired cycles. Code in R, SAS, and GenStat
is provided. The paper clarifies that the appropriate sample
size criterion depends fundamentally on whether the goal is
individual-level inference, population-level inference, or
detection of treatment-response heterogeneity.

\medskip
\textbf{Schmid CH, Duan N, Kravitz RL.} Statistical design
and analytic considerations for N-of-1 trials. Agency for
Healthcare Research and Quality. 2014.

\smallskip\noindent
This AHRQ technical report provides a comprehensive
treatment of statistical methods for N-of-1 trial design.
The authors cover paired-cycle randomization schemes,
sample size calculations for both individual and aggregated
trials, and analytic approaches including mixed-effects
models and Bayesian methods. Design factors addressed
include the number of cycles, washout period duration, and
measurement frequency. The report serves as a practical
reference for investigators planning simulation studies of
N-of-1 designs.

\medskip
\textbf{Pequignat F, Balaam B, Sherrington R, Sherrington
C, Maher C.} Power analysis for idiographic
(within-subject) clinical trials: implications for
treatments of rare conditions and precision medicine.
\textit{Behavior Research Methods}. 2023;55:3824--3842.
doi:10.3758/s13428-022-02012-1.

\smallskip\noindent
Pequignat et al.\ present simulation-based power analysis
for within-subject clinical trials, with particular
attention to rare conditions where aggregated designs are
infeasible. The simulation framework varies effect size,
autocorrelation, number of measurement occasions, and
number of treatment cycles. The findings quantify how
autocorrelation inflates type~I error when ignored and how
the required number of measurements increases rapidly for
small effect sizes, offering practical guidance for
precision medicine applications.

\medskip
\textbf{Senn S, Julious S, Chen Y.} A methodological
review of randomised N-of-1 trials. \textit{Trials}.
2024;25:264. doi:10.1186/s13063-024-08100-1.

\smallskip\noindent
This methodological review examines the design and
analysis characteristics of published randomized N-of-1
trials. Senn, Julious, and Chen assess how commonly
investigators address key design elements such as
randomization schemes, washout periods, blinding, and
outcome measurement frequency. The review identifies
widespread inconsistencies in reporting and analytic
practice, providing an empirical basis for simulation
studies that compare the consequences of different design
choices.

% --------------------------------------------------------
\section{Carryover and Temporal Correlation}
% --------------------------------------------------------

\textbf{Yap C, Wang X, Harhay MO.} Behavioral carry-over
effect and power consideration in crossover trials.
\textit{Biometrics}. 2024;80(2):ujae023.
doi:10.1093/biomtcs/ujae023.

\smallskip\noindent
Yap, Wang, and Harhay develop a formal framework for
behavioral carryover effects in crossover trials, where
carryover operates through psychological or behavioral
mechanisms rather than pharmacokinetic ones. They use
simulation to evaluate power under various carryover
magnitudes and asymmetries, showing that behavioral
carryover can bias treatment effect estimates even when
pharmacokinetic carryover is negligible. The work is
directly relevant to prazosin-PTSD trials, where
psychological expectancy effects may persist across
treatment periods.

\medskip
\textbf{Sturdevant SG, Lumley T.} Statistical methods for
testing carryover effects: a mixed effects model approach.
\textit{Contemporary Clinical Trials Communications}.
2021;22:100711. doi:10.1016/j.conctc.2021.100711.

\smallskip\noindent
Sturdevant and Lumley propose a mixed-effects model
approach for testing carryover effects and validate it
through simulation based on blood pressure measurements.
The simulations demonstrate that modeling the continuous
outcome directly, rather than dichotomized thresholds,
yields valid tests even when open-label treatment after
diagnosis censors subsequent observations. The key insight
is that the censored data are missing at random under the
proposed model, enabling consistent estimation of carryover
parameters.

\medskip
\textbf{Schork NJ.} Accommodating serial correlation and
sequential design elements in personalized studies and
aggregated personalized studies. \textit{Harvard Data
Science Review}. 2022;SI3.
doi:10.1162/99608f92.f1eef6f4.

\smallskip\noindent
Schork investigates the consequences of serial correlation
for the design and analysis of personalized studies through
extensive simulation. The paper explores sequential
stopping rules that allow real-time decisions about
treatment efficacy, the frequency of measurement collection,
and the interaction between washout periods and
autocorrelation structure. Simulation results show that
ignoring serial correlation can substantially inflate false
positive rates and that sequential designs offer efficiency
gains when correlation is properly modeled.

\medskip
\textbf{Tang W, Landes RD.} Some t-tests for N-of-1 trials
with serial correlation. \textit{PLoS One}.
2020;15(2):e0228077. doi:10.1371/journal.pone.0228077.

\smallskip\noindent
Tang and Landes develop modified t-tests that account for
serial correlation in N-of-1 trial data and evaluate their
operating characteristics via simulation. The simulation
study varies the autocorrelation coefficient, the number of
observations per period, and the treatment effect size.
The proposed tests maintain nominal type~I error rates
under AR(1) correlation structures where standard t-tests
fail, demonstrating the practical importance of correlation
adjustment in single-subject analyses.

\medskip
\textbf{arXiv preprint.} Estimation of complex carryover
effects in crossover designs with repeated measures.
arXiv:2402.16362. 2024.

\smallskip\noindent
This preprint addresses the estimation of complex carryover
patterns in crossover designs with repeated measurements
within each period. The authors use simulation to compare
estimators under multi-order and asymmetric carryover
structures that go beyond the simple first-order models
commonly assumed. The results are relevant to N-of-1 trial
simulation because they characterize bias and efficiency
loss when carryover models are misspecified.

% --------------------------------------------------------
\section{Biomarker-Treatment Interactions}
% --------------------------------------------------------

\textbf{Haller B, Ulm K.} A simulation study on estimating
biomarker-treatment interaction effects in randomized
trials with prognostic variables. \textit{Trials}.
2018;19:128. doi:10.1186/s13063-018-2491-0.

\smallskip\noindent
Haller and Ulm conduct a simulation study examining how
prognostic variables affect the estimation of
biomarker-treatment interactions in randomized trials with
time-to-event outcomes. They compare strategies of ignoring
prognostic factors, including all covariates, and applying
variable selection. The simulations reveal that omitting
relevant prognostic variables biases the interaction
estimate, while including them increases detection power.
However, indiscriminate inclusion or poor variable selection
elevates false positive rates.

\medskip
\textbf{Grenet G, Blanc C, Bardel C, Gueyffier F, Roy P.}
Comparison of crossover and parallel-group designs for the
identification of a binary predictive biomarker of the
treatment effect. \textit{Basic \& Clinical Pharmacology \&
Toxicology}. 2020;126:59--64. doi:10.1111/bcpt.13293.

\smallskip\noindent
Grenet et al.\ use simulation to compare crossover and
parallel-group designs for identifying a binary predictive
biomarker. The simulation framework varies sample size,
biomarker prevalence, interaction effect magnitude, and
within-subject correlation. Results indicate that crossover
designs are more efficient for detecting predictive
biomarkers when within-subject correlation is high, a
finding that directly supports the use of N-of-1 designs
for biomarker validation in the pmsimstats framework.

% --------------------------------------------------------
\section{Analysis Methods Evaluated via Simulation}
% --------------------------------------------------------

\textbf{Senn S, Julious SA, Araujo A.} A comparison of four
methods for the analysis of N-of-1 trials. \textit{PLoS
One}. 2014;9(2):e87752. doi:10.1371/journal.pone.0087752.

\smallskip\noindent
Senn, Julious, and Araujo compare four analytic methods
for N-of-1 trial data (fixed-effects, random-effects,
Bayesian hierarchical, and a simple paired analysis)
through simulation. The simulation design varies the number
of patients, the number of cycles, the magnitude of
treatment-by-patient interaction variance, and the
within-patient variance. The comparison reveals that all
methods perform similarly for population-level inference
but differ substantially in their ability to estimate
individual treatment effects, with hierarchical models
offering the best shrinkage properties.

\medskip
\textbf{Senn S, Araujo A, Julious S.} Understanding
variation in sets of N-of-1 trials. \textit{PLoS One}.
2016;11(12):e0167167.

\smallskip\noindent
This paper uses both analytical derivations and simulation
to investigate how variance components in sets of N-of-1
trials affect inference about individual treatment
responses. The authors decompose observed variation into
within-patient and between-patient components and show
through simulation that the commonly held belief in strong
personal treatment response elements is often not supported
by the data. The work has implications for simulation study
design, as it clarifies which variance structures are
realistic for generating synthetic trial data.

\medskip
\textbf{Chen Y, Li P, Yang K.} Comparison of aggregated
N-of-1 trials with parallel and crossover randomized
controlled trials using simulation studies. \textit{PLoS
One}. 2020;15(1):e0216949.
doi:10.1371/journal.pone.0216949.

\smallskip\noindent
Chen, Li, and Yang replicate and extend the Blackston et
al.\ simulation framework, comparing aggregated N-of-1,
parallel, and crossover designs across additional variance
structures. The simulation systematically varies patient
heterogeneity, error variance, and carryover magnitude.
Their results corroborate the power advantages of N-of-1
designs and further quantify the conditions under which
carryover bias becomes problematic, providing independent
validation of the Blackston findings.

\medskip
\textbf{Senn S, Julious S.} The analysis of continuous data
from N-of-1 trials using paired cycles: a simple tutorial.
\textit{Trials}. 2024;25:123.
doi:10.1186/s13063-024-07964-7.

\smallskip\noindent
Senn and Julious present a tutorial on paired-cycle
analysis for continuous N-of-1 trial data, illustrating
methods through worked examples. Although primarily
pedagogical, the paper includes simulation-based
verification that the paired-cycle approach maintains
appropriate coverage and type~I error control. The tutorial
clarifies the relationship between design balance,
randomization constraints, and the validity of simple
paired analyses.

\medskip
\textbf{Peterlin J, Kej\v{z}ar N, Blagus R.} Correct
specification of design matrices in linear mixed effects
models: tests with graphical representation. \textit{TEST}.
2023;32:184--210. doi:10.1007/s11749-022-00830-1.

\smallskip\noindent
Peterlin, Kej\v{z}ar, and Blagus develop diagnostic tests
for verifying the correct specification of fixed- and
random-effects design matrices in linear mixed models,
validated through a large Monte Carlo simulation study.
The simulation encompasses multilevel and crossed random
effects structures relevant to N-of-1 trial analysis. The
proposed graphical diagnostics enable researchers to detect
misspecification that would otherwise compromise the
validity and efficiency of treatment effect estimates.

% --------------------------------------------------------
\section{Bayesian Approaches}
% --------------------------------------------------------

\textbf{Schmid CH, Staudenmayer J.} Bayesian models for
N-of-1 trials. \textit{Harvard Data Science Review}. 2021.
doi:10.1162/99608f92.c06f5e3d.

\smallskip\noindent
Schmid and Staudenmayer review Bayesian modeling approaches
for N-of-1 trials, covering hierarchical models that borrow
strength across patients and time periods. The authors
discuss prior specification, posterior inference for
individual treatment effects, and the role of simulation in
calibrating Bayesian models. The paper provides a
conceptual foundation for simulation studies that compare
frequentist and Bayesian analytic strategies in the
aggregated N-of-1 setting.

\medskip
\textbf{De Bock M, Van den Noortgate W, Gielen M.} Accurate
models vs.\ accurate estimates: a simulation study of
Bayesian single-case experimental designs. \textit{Behavior
Research Methods}. 2021;53:1782--1798.
doi:10.3758/s13428-020-01522-0.

\smallskip\noindent
De Bock, Van den Noortgate, and Gielen use simulation to
investigate the tradeoff between model complexity and
estimation accuracy in Bayesian single-case designs. The
simulation varies the number of measurement occasions, the
autocorrelation structure, and the complexity of the
mean model. Results indicate that simpler models sometimes
yield more accurate treatment effect estimates despite
being technically misspecified, a finding relevant to
choosing analysis models in N-of-1 simulation frameworks.

\medskip
\textbf{Daza EJ, Wac K, Oppezzo M.} Causal analysis of
self-tracked time series data using a counterfactual
framework for N-of-1 trials. \textit{AMIA Annual Symposium
Proceedings}. 2018;2018:532--541.

\smallskip\noindent
Daza, Wac, and Oppezzo propose a counterfactual framework
for causal inference in N-of-1 trials using self-tracked
time series data. The approach employs simulation of
potential outcomes under treatment and control to estimate
individual causal effects. The framework addresses the
challenge of time-varying confounding in observational
N-of-1 data and provides a bridge between the traditional
randomized N-of-1 design and emerging digital health
applications.

% --------------------------------------------------------
\section{Clinical Context: Prazosin-PTSD Trials}
% --------------------------------------------------------

The prazosin-PTSD trial program provides the motivating
clinical application for the pmsimstats simulation
framework. The progression from small proof-of-concept
studies to a large VA cooperative trial, followed by an
unexpected null result, illustrates the statistical
challenges that simulation methods can help address.

\medskip
\textbf{Raskind MA, Peskind ER, Kanter ED, et al.}
Reduction of nightmares and other PTSD symptoms in combat
veterans by prazosin: a placebo-controlled study.
\textit{American Journal of Psychiatry}.
2003;160(2):371--373.

\smallskip\noindent
This initial proof-of-concept study (N=10) demonstrated
significant reduction in trauma nightmares and overall PTSD
symptom severity with prazosin versus placebo in a
double-blind crossover design. The small sample and
crossover structure make this an early example of the
design considerations that motivate simulation-based power
analysis for N-of-1 trials in the PTSD context.

\medskip
\textbf{Raskind M, Peskind E, McFall M, Peterson K, Engel
C, Doyle M.} A placebo-controlled trial of prazosin vs.\
paroxetine in combat stress-induced PTSD nightmares and
sleep disturbance. Department of Defense Technical Report.
2007. doi:10.21236/ada484244.

\smallskip\noindent
This report describes a comparative trial of prazosin and
paroxetine for combat-related PTSD nightmares. The study
design required careful consideration of differential
carryover between the two active agents. The results
supported prazosin efficacy and highlighted the challenge of
modeling pharmacologically distinct carryover profiles in
crossover designs.

\medskip
\textbf{Raskind M, Peskind E, Peterson K, Engel C.} A
placebo-controlled augmentation trial of prazosin for
combat trauma PTSD. Department of Defense Technical Report.
2009. doi:10.21236/ada515114.

\smallskip\noindent
This augmentation trial evaluated prazosin added to
existing treatment in a parallel-group design. The shift
from crossover to parallel-group architecture reflects the
practical challenges of washout periods and carryover in
PTSD pharmacotherapy, providing empirical context for
simulation studies that compare design efficiency across
trial architectures.

\medskip
\textbf{Raskind MA, Peskind ER, Williams T, et al.} A trial
of prazosin for combat trauma PTSD with nightmares in
active-duty soldiers returned from Iraq and Afghanistan.
\textit{American Journal of Psychiatry}.
2013;170(9):1003--1010.
doi:10.1176/appi.ajp.2013.12081133.

\smallskip\noindent
This parallel-group RCT (N=67) in active-duty soldiers
found significant prazosin effects on nightmares and sleep
quality. The moderate sample size and observed effect sizes
provide empirical parameters for calibrating simulation
studies. The study also illustrates the heterogeneity of
treatment response across subpopulations, motivating
biomarker-stratified design approaches.

\medskip
\textbf{Raskind MA, Peskind ER, Hart KL, et al.} Trial of
prazosin for post-traumatic stress disorder in military
veterans. \textit{New England Journal of Medicine}.
2018;378(6):507--517. doi:10.1056/NEJMoa1507598.

\smallskip\noindent
The large VA cooperative study (N=304) found no significant
difference between prazosin and placebo on the primary
outcome, contradicting prior smaller trials. This
unexpected null result is the central motivating puzzle for
the pmsimstats simulation framework. The discrepancy
raises questions about population heterogeneity, biomarker
moderation, site effects, and the adequacy of
parallel-group designs for detecting effects that may be
concentrated in biomarker-defined subgroups.

\medskip
\textbf{Raskind MA, Peskind ER, Poupore EL, et al.}
Randomized controlled pilot trial of prazosin for
prophylaxis of posttraumatic headaches in active-duty
service members and veterans. \textit{Headache}.
2023;63(6):751--762. doi:10.1111/head.14529.

\smallskip\noindent
This pilot study extends the prazosin program to
posttraumatic headache prophylaxis using a parallel-group
design. The small sample (N=48) and exploratory nature
exemplify the conditions under which N-of-1 or aggregated
crossover designs could provide more efficient evidence
generation. The study parameters inform simulation
scenarios for related neuropsychiatric outcomes.

\vfill
\noindent\rule{\textwidth}{0.4pt}
{\footnotesize Rendered on 2026-04-09 at 18:49 PDT.\\
Source: \textasciitilde/prj/alz/10-pmsimstats-ng/pmsimstats-ng/docs/20-annotated-bibliography.tex}

\end{document}
